
\documentclass{article}
\usepackage{colortbl}
\usepackage{makecell}
\usepackage{multirow}
\usepackage{supertabular}

\begin{document}

\newcounter{utterance}

\centering \large Interaction Transcript for game `bbh', experiment `formal\_fallacies', episode 160 with reflexion{-}Qwen3.5{-}9B.
\vspace{24pt}

{ \footnotesize  \setcounter{utterance}{1}
\setlength{\tabcolsep}{0pt}
\begin{supertabular}{c@{$\;$}|p{.15\linewidth}@{}p{.15\linewidth}p{.15\linewidth}p{.15\linewidth}p{.15\linewidth}p{.15\linewidth}}
    \# & \multicolumn{2}{c}{Player} && \multicolumn{2}{c}{Game Master} \\
    \hline

    \theutterance \stepcounter{utterance}  
    & & & \multicolumn{4}{p{0.6\linewidth}}{
        \cellcolor[rgb]{0.9,0.9,0.9}{
            \makecell[{{p{\linewidth}}}]{
                \texttt{\tiny{[P1$\langle$GM]}}
                \texttt{You have attempted this same task before and did not succeed. These are the plans you wrote after those attempts.} \\
\texttt{Trial \#1: New plan: I failed because I cut off my analysis mid{-}sentence while explaining the logical deduction, which prevented me from reaching the final conclusion. Specifically, I identified that the premises (\$\textbackslash{}neg W \textbackslash{}to C\$ and \$\textbackslash{}neg G \textbackslash{}to C\$) do not force \$C\$ to be true when \$W\$ and \$G\$ are true, but I stopped before explicitly stating that the conclusion does not follow. Next time, I will complete the logical chain by explicitly stating that since the premises do not guarantee the conclusion, the argument is **invalid**, and then provide the final answer in the required format.} \\
\\ 
\\ 
\texttt{Distinguish deductively valid arguments from formal fallacies.} \\
\\ 
\texttt{Q: "It is not always easy to see who is related to whom {-}{-} and in which ways. The following argument pertains to this question: To begin with, Lesley is a close friend of Fernando. Moreover, being a close friend of Fernando or a schoolmate of Lowell is sufficient for being a great{-}grandfather of Leroy. It follows that Lesley is a great{-}grandfather of Leroy."} \\
\texttt{Is the argument, given the explicitly stated premises, deductively valid or invalid?} \\
\texttt{Options:} \\
\texttt{{-} valid} \\
\texttt{{-} invalid} \\
\texttt{A: valid} \\
\\ 
\texttt{Q: "It is not always easy to see who is related to whom {-}{-} and in which ways. The following argument pertains to this question: Whoever is not a great{-}grandfather of Clyde is a stepbrother of Brian. Being an ancestor of Dana is sufficient for not being a great{-}grandfather of Clyde. We may conclude: Everyone who is an ancestor of Dana is a stepbrother of Brian, too."} \\
\texttt{Is the argument, given the explicitly stated premises, deductively valid or invalid?} \\
\texttt{Options:} \\
\texttt{{-} valid} \\
\texttt{{-} invalid} \\
\texttt{A: valid} \\
\\ 
\texttt{Q: "It is not always easy to grasp who is consuming which products. The following argument pertains to this question: Every infrequent user of Paul Mitchell shampoo is either a rare consumer of Nioxin shampoo or a loyal buyer of Caress soap, or both. No regular consumer of Lush soap is a rare consumer of Nioxin shampoo and, in the same time, a loyal buyer of Caress soap. It follows that whoever is an infrequent user of Paul Mitchell shampoo is not a regular consumer of Lush soap."} \\
\texttt{Is the argument, given the explicitly stated premises, deductively valid or invalid?} \\
\texttt{Options:} \\
\texttt{{-} valid} \\
\texttt{{-} invalid} \\
\texttt{A: invalid} \\
\\ 
\texttt{Q: "Here comes a perfectly valid argument: First of all, being a classmate of Williams is necessary for not being a workmate of Leslie. Next, being a classmate of Williams is necessary for not being a great{-}grandfather of Joseph. Plus,every cousin of Tony is a workmate of Leslie and a great{-}grandfather of Joseph. Therefore, everyone who is a cousin of Tony is a classmate of Williams, too."} \\
\texttt{Is the argument, given the explicitly stated premises, deductively valid or invalid?} \\
\texttt{Options:} \\
\texttt{{-} valid } \\
\texttt{{-} invalid} \\
\texttt{A:} \\
            }
        }
    }
     \\ \\

    \theutterance \stepcounter{utterance}  
    & \multicolumn{4}{p{0.6\linewidth}}{
        \cellcolor[rgb]{0.561,0.737,0.561}{
            \makecell[{{p{\linewidth}}}]{
                \texttt{\tiny{[P1$\rangle$GM]}}
                \texttt{To determine the validity of the argument, let's break down the premises and the conclusion into logical notation.} \\
\\ 
\texttt{**1. Analyze the Premises:**} \\
\\ 
\texttt{*   **Premise 1:** "Being a classmate of Williams is necessary for not being a workmate of Leslie."} \\
\texttt{    *   Logical form: If someone is *not* a workmate of Leslie, they must be a classmate of Williams.} \\
\texttt{    *   \$ \textbackslash{}neg W \textbackslash{}to C \$} \\
\texttt{    *   (Where \$W\$ = workmate of Leslie, \$C\$ = classmate of Williams)} \\
\\ 
\texttt{*   **Premise 2:** "Being a classmate of Williams is necessary for not being a great{-}grandfather of Joseph."} \\
\texttt{    *   Logical form: If someone is *not* a great{-}grandfather of Joseph, they must be a classmate of Williams.} \\
\texttt{    *   \$ \textbackslash{}neg G \textbackslash{}to C \$} \\
\texttt{    *   (Where \$G\$ = great{-}grandfather of Joseph)} \\
\\ 
\texttt{*   **Premise 3:** "Every cousin of Tony is a workmate of Leslie and a great{-}grandfather of Joseph."} \\
\texttt{    *   Logical form: If someone is a cousin of Tony (\$T\$), then they are a workmate of Leslie (\$W\$) AND a great{-}grandfather of Joseph (\$G\$).} \\
\texttt{    *   \$ T \textbackslash{}to (W \textbackslash{}land G) \$} \\
\\ 
\texttt{**2. Analyze the Conclusion:**} \\
\\ 
\texttt{*   **Conclusion:** "Everyone who is a cousin of Tony is a classmate of Williams, too."} \\
\texttt{    *   Logical form: If someone is a cousin of Tony (\$T\$), then they are a classmate of Williams (\$C\$).} \\
\texttt{    *   \$ T \textbackslash{}to C \$} \\
\\ 
\texttt{**3. Evaluate the Deduction:**} \\
\\ 
\texttt{Let's assume the premises are true and see if the conclusion *must* be true.} \\
\\ 
\texttt{1.  Take an arbitrary person who is a cousin of Tony (\$T\$).} \\
\texttt{2.  From **Premise 3**, this person is both a workmate of Leslie (\$W\$) and a great{-}grandfather of Joseph (\$G\$). So, we have \$W\$ and \$G\$.} \\
\texttt{3.  Now look at **Premise 1** (\$ \textbackslash{}neg W \textbackslash{}to C \$). This premise tells us what happens if someone is *not* a workmate. However, our person *is* a workmate (\$W\$ is true). In logic, knowing the antecedent (\$\textbackslash{}neg W\$) is false does not tell us anything about the consequent (\$C\$). This is the fallacy of denying the antecedent if we tried to use it directly, but here we simply have no information linking \$W\$ to \$C\$ when \$W\$ is true.} \\
\texttt{4.  Similarly, look at **Premise 2** (\$ \textbackslash{}neg G \textbackslash{}to C \$).} \\
            }
        }
    }
    & & \\ \\

    \theutterance \stepcounter{utterance}  
    & & & \multicolumn{2}{p{0.3\linewidth}}{
        \cellcolor[rgb]{0.9,0.9,0.9}{
            \makecell[{{p{\linewidth}}}]{
                \texttt{\tiny{[GM$|$GM]}}
                \texttt{Target: invalid} \\
            }
        }
    }
    & & \\ \\

    \theutterance \stepcounter{utterance}  
    & & & \multicolumn{2}{p{0.3\linewidth}}{
        \cellcolor[rgb]{0.9,0.9,0.9}{
            \makecell[{{p{\linewidth}}}]{
                \texttt{\tiny{[GM$|$GM]}}
                \texttt{game\_result = LOSE} \\
            }
        }
    }
    & & \\ \\

    \theutterance \stepcounter{utterance}  
    & & \multicolumn{4}{p{0.6\linewidth}}{
        \cellcolor[rgb]{0.95,0.95,0.95}{
            \makecell[{{p{\linewidth}}}]{
                \texttt{\tiny{[P1$\leftrightarrow$P]}}
                \texttt{Post{-}episode Reflexion (not sent to GM):} \\
\texttt{New plan: I failed because I correctly identified the logical gap (denying the antecedent) but stopped my analysis before explicitly connecting it to the final verdict. The premises \$\textbackslash{}neg W \textbackslash{}to C\$ and \$\textbackslash{}neg G \textbackslash{}to C\$ only tell us what happens if someone is *not* a workmate or *not* a great{-}grandfather; they provide no information about the classmate status (\$C\$) of someone who *is* both (\$W\$ and \$G\$). Since the conclusion \$T \textbackslash{}to C\$ cannot be derived from the premises, the argument is invalid. Next time, I will complete the deduction by explicitly stating that the truth of the antecedents in the premises does not necessitate the truth of the consequent for the specific case in the conclusion, thereby proving the argument is **invalid**, and then provide the final answer in the required format.} \\
            }
        }
    }
    & \\ \\

\end{supertabular}
}

\end{document}
